\documentclass[a4paper,12pt]{article}

\usepackage{fontspec}
\setmainfont{PT Serif}
\usepackage{babel}
\babelprovide[import,main]{russian}
\usepackage{microtype}
\usepackage[margin=2.5cm]{geometry}

\setlength{\parindent}{1.25cm}
\setlength{\parskip}{0pt}

\begin{document}

\begin{center}
    {\Large\bfseries Теория вероятностей}

    \vspace{0.3em}

    {\large Лекция 4}

    \vspace{0.3em}

    \textbf{Условная вероятность}

    \vspace{0.3em}

    {\small 12 сентября 2026 года}
\end{center}

\vspace{1em}

\section{Условная вероятность}

Текст лекции раскрывает тему. Порядок взят из исходного материала.

% Здесь могут находиться любые подходящие bricks.

\subsection{Независимость событий}

Следующий смысловой блок лекции. Подраздел используется только тогда, когда
такая иерархия действительно присутствует в исходнике.

\section{Формула полной вероятности}

Продолжение теоретического материала.

\end{document}
